\documentclass[../DoAn.tex]{subfiles}
\begin{document}

Từ các yêu cầu chức năng và phi chức năng đã xác định ở Chương~\ref{chapter:Related_works}, chương này trình bày những nền tảng lý thuyết và công nghệ cốt lõi được sử dụng để xây dựng hệ thống. Với mỗi công nghệ, đồ án nêu rõ vấn đề mà nó giải quyết, một số lựa chọn thay thế đã cân nhắc và lý do lựa chọn, nhằm bảo đảm tính nhất quán với các yêu cầu đã đặt ra ở chương trước.

\section{Nền tảng lý thuyết}
\label{section:3.1}
Hệ thống dựa trên một số nguyên lý đã được kiểm chứng trong khoa học học tập. Học tập chủ động cho thấy người học ghi nhớ và vận dụng kiến thức tốt hơn khi tham gia trực tiếp thay vì tiếp nhận một chiều~\cite{prince2004active}, còn hiệu ứng kiểm tra chỉ ra rằng việc truy hồi kiến thức qua câu hỏi giúp củng cố trí nhớ và tránh ảo giác đã hiểu bài~\cite{karpicke2008critical}. Với nội dung truyền tải qua video, các nguyên lý thiết kế video giáo dục và phân đoạn nội dung khuyến nghị chia bài giảng thành những đoạn ngắn kèm yếu tố tương tác để giảm tải nhận thức~\cite{brame2016effective,mayer2003nine}; bên cạnh đó, phản hồi tức thời sau mỗi hoạt động là một trong những yếu tố ảnh hưởng mạnh nhất tới kết quả học tập~\cite{hattie2007power}. Những nguyên lý này là cơ sở cho định hướng cốt lõi của hệ thống: chèn hoạt động tương tác vào từng mốc thời gian của video, buộc người học truy hồi kiến thức và nhận phản hồi ngay thay vì xem thụ động.

\section{Công nghệ xử lý trí tuệ nhân tạo}
\label{section:3.2}
Khối xử lý trí tuệ nhân tạo là thành phần cốt lõi giúp hệ thống tự động chuyển một video bài giảng thành bài học tương tác. Khối này kết hợp ba loại công nghệ: nhận dạng tiếng nói để lấy lời thoại, mô hình ngôn ngữ lớn để phân đoạn nội dung và sinh hoạt động, và tổng hợp tiếng nói để tạo bài nghe. Toàn bộ quy trình được minh họa ở Hình~\ref{fig:pipeline_ai}: video bài giảng lần lượt đi qua các giai đoạn phiên âm, phân đoạn nội dung và sinh hoạt động tương tác, toàn bộ chạy ở chế độ nền.

\begin{figure}[H]
\centering
\includegraphics[width=0.95\textwidth]{Hinhve/chuong_3/quytrinhai.pdf}
\caption{Quy trình xử lý trí tuệ nhân tạo sinh bài học}
\label{fig:pipeline_ai}
\end{figure}

\subsection{Nhận dạng tiếng nói}
Để phân tích nội dung bài giảng, hệ thống cần chuyển lời thoại trong video thành văn bản, đáp ứng yêu cầu trích xuất bản chép lời ở Chương~\ref{chapter:Related_works}. Nhiệm vụ này do mô hình nhận dạng tiếng nói Whisper~\cite{radford2023whisper} đảm nhận, được triển khai tại chỗ trên hạ tầng riêng. Whisper có một số đặc điểm phù hợp với bài toán:

\textbf{Độ chính xác đa ngôn ngữ} -- Whisper là mô hình mã nguồn mở được huấn luyện trên khối dữ liệu âm thanh quy mô lớn, cho độ chính xác cao và hỗ trợ tốt tiếng Việt.

\textbf{Mốc thời gian} -- kết quả nhận dạng kèm theo mốc thời gian, cần thiết để gắn các hoạt động tương tác vào đúng vị trí trong video.

\textbf{Triển khai tại chỗ} -- mô hình chạy trên hạ tầng riêng nên dữ liệu bài giảng không phải gửi ra dịch vụ bên thứ ba.

Các lựa chọn thay thế gồm dịch vụ đám mây như Google Cloud Speech-to-Text hay Amazon Transcribe, và thư viện ngoại tuyến nhẹ như Vosk. Bảng~\ref{table:stt} so sánh các lựa chọn theo độ chính xác, chi phí, khả năng giữ dữ liệu tại chỗ và công sức vận hành.

\begin{table}[ht]
\centering
\small
\caption{So sánh các công nghệ nhận dạng tiếng nói}
\label{table:stt}
\begin{tabular}{|>{\raggedright\arraybackslash}p{3cm}|>{\raggedright\arraybackslash}p{3.2cm}|>{\raggedright\arraybackslash}p{3.7cm}|>{\raggedright\arraybackslash}p{2.6cm}|}
\hline
\textbf{Tiêu chí} & \textbf{Whisper (tự triển khai)} & \textbf{Dịch vụ đám mây (Google, Amazon)} & \textbf{Vosk} \\ \hline
Độ chính xác tiếng Việt & Cao & Cao & Trung bình \\ \hline
Chi phí & Miễn phí & Trả theo lượt gọi & Miễn phí \\ \hline
Dữ liệu giữ tại chỗ & Có & Không & Có \\ \hline
Mốc thời gian chi tiết & Có & Có & Hạn chế \\ \hline
Công sức vận hành & Trung bình & Thấp & Thấp \\ \hline
\end{tabular}
\end{table}

Như Bảng~\ref{table:stt} cho thấy, Whisper tự triển khai là phương án duy nhất vừa đạt độ chính xác cao, vừa miễn phí và giữ được dữ liệu bài giảng tại chỗ, nên được lựa chọn.

\subsection{Mô hình ngôn ngữ lớn}
Sau khi có bản chép lời, hệ thống tự động phân đoạn nội dung và sinh ra các hoạt động tương tác, đây là yêu cầu trung tâm của đề tài ở Chương~\ref{chapter:Related_works}. Nhiệm vụ này được giao cho mô hình ngôn ngữ lớn, cụ thể là Gemini~\cite{geminiteam2023gemini} của Google. Các mô hình ngôn ngữ lớn hiện đại dựa trên kiến trúc Transformer~\cite{vaswani2017attention} và có những đặc điểm sau:

\textbf{Sinh nội dung có cấu trúc} -- mô hình tạo câu hỏi và hoạt động ở định dạng có cấu trúc, nhờ đó hệ thống tự động đọc và lưu lại được kết quả sinh ra (Hình~\ref{fig:gemini}).

\textbf{Hiểu ngữ cảnh dài} -- mô hình nắm bắt được toàn bộ một đoạn nội dung để đặt câu hỏi bám sát bài giảng và chấm tự động các câu trả lời tự luận.

\textbf{Khả năng thay thế} -- thiết kế cho phép dễ dàng đổi sang một nhà cung cấp mô hình khác khi cần, không phụ thuộc cứng vào một dịch vụ.

\begin{figure}[H]
\centering
\includegraphics[width=0.9\textwidth]{Hinhve/chuong_3/gemini.pdf}
\caption{Mô hình ngôn ngữ lớn sinh hoạt động tương tác có cấu trúc từ một đoạn nội dung}
\label{fig:gemini}
\end{figure}

Các lựa chọn thay thế gồm GPT của OpenAI và các mô hình mã nguồn mở tự triển khai như Llama hay Qwen. Bảng~\ref{table:llm} so sánh các lựa chọn theo chất lượng đầu ra có cấu trúc, năng lực tiếng Việt, chi phí và khả năng chạy tại chỗ.

\begin{table}[ht]
\centering
\small
\caption{So sánh các mô hình ngôn ngữ lớn}
\label{table:llm}
\begin{tabular}{|>{\raggedright\arraybackslash}p{2.9cm}|>{\raggedright\arraybackslash}p{3cm}|>{\raggedright\arraybackslash}p{3cm}|>{\raggedright\arraybackslash}p{3.4cm}|}
\hline
\textbf{Tiêu chí} & \textbf{Gemini} & \textbf{GPT (OpenAI)} & \textbf{Llama, Qwen (tự triển khai)} \\ \hline
Đầu ra có cấu trúc & Cao & Cao & Trung bình \\ \hline
Năng lực tiếng Việt & Cao & Cao & Trung bình \\ \hline
Chi phí & Có gói miễn phí & Trả phí & Miễn phí (cần phần cứng) \\ \hline
Chạy tại chỗ & Không & Không & Có \\ \hline
\end{tabular}
\end{table}

Như Bảng~\ref{table:llm} cho thấy, Gemini đạt chất lượng đầu ra và năng lực tiếng Việt ở mức cao tương đương GPT nhưng có gói miễn phí phục vụ phát triển, nên được lựa chọn.

\subsection{Tổng hợp tiếng nói}
Một số hoạt động, đặc biệt là bài nghe ở Chương~\ref{chapter:Related_works}, cần chuyển văn bản thành âm thanh. Hệ thống dùng Piper~\cite{piper2023}, một công cụ tổng hợp tiếng nói bằng mạng nơ-ron, với hai đặc điểm phù hợp với bài toán:

\textbf{Chạy trên CPU} -- Piper là mô hình mã nguồn mở, hoạt động hoàn toàn trên CPU thông thường nên không đòi hỏi phần cứng chuyên dụng.

\textbf{Hỗ trợ tiếng Việt} -- Piper có sẵn giọng đọc tiếng Việt, phù hợp để tạo nội dung bài nghe cho người học.

Các lựa chọn thay thế gồm dịch vụ đám mây như Google Cloud Text-to-Speech hay Azure (chất lượng cao nhưng trả phí và phải gửi dữ liệu ra ngoài) và mô hình VieNeu-TTS (giọng tự nhiên hơn nhưng đòi hỏi GPU). Piper được ưu tiên vì chạy trên CPU thông thường, không phát sinh chi phí và dễ triển khai cùng hạ tầng sẵn có.

\section{Công nghệ phía máy chủ}
\label{section:3.3}
Phía máy chủ chịu trách nhiệm xử lý nghiệp vụ, lưu trữ dữ liệu và điều phối các tác vụ phân tích chạy nền. Hình~\ref{fig:kientruc} trình bày kiến trúc tổng thể của hệ thống: người dùng truy cập qua trình duyệt, một máy chủ proxy đứng trước điều hướng yêu cầu tới giao diện web hoặc máy chủ ứng dụng; máy chủ ứng dụng kết nối tới các kho lưu trữ và các dịch vụ trí tuệ nhân tạo đã trình bày ở Phần~\ref{section:3.2}. Phần này trình bày các công nghệ chính ở phía máy chủ cùng lý do lựa chọn.

\begin{figure}[H]
\centering
\includegraphics[width=0.95\textwidth]{Hinhve/chuong_3/kientructongthe.pdf}
\caption{Kiến trúc tổng thể của hệ thống}
\label{fig:kientruc}
\end{figure}

\subsection{Ngôn ngữ lập trình Go}
Máy chủ được xây dựng bằng ngôn ngữ Go~\cite{golang}, đáp ứng các yêu cầu phi chức năng về hiệu năng và khả năng mở rộng ở Chương~\ref{chapter:Related_works}.

\textbf{Xử lý đồng thời} -- Go cung cấp cơ chế goroutine nhẹ, thuận lợi cho việc chạy song song nhiều tác vụ phân tích và phục vụ nhiều yêu cầu cùng lúc (Hình~\ref{fig:goworker}).

\begin{figure}[H]
\centering
\includegraphics[width=0.5\textwidth]{Hinhve/chuong_3/goworker.pdf}
\caption{Nhiều worker (goroutine) lấy việc từ hàng đợi và xử lý song song}
\label{fig:goworker}
\end{figure}

\textbf{Hiệu năng} -- chương trình được biên dịch thành mã máy chạy nhanh và ít tiêu tốn tài nguyên.

\textbf{Dễ bảo trì} -- Go có cú pháp đơn giản và bộ công cụ chuẩn hóa, giúp mã nguồn dễ đọc và dễ bảo trì.

Các lựa chọn thay thế gồm Node.js (dùng chung ngôn ngữ với giao diện nhưng yếu hơn ở các tác vụ nặng về tính toán) và Java cùng Spring (hệ sinh thái trưởng thành nhưng môi trường chạy nặng hơn). Go được chọn nhờ cân bằng tốt giữa hiệu năng, khả năng xử lý đồng thời và độ gọn nhẹ.

\subsection{Giao tiếp Connect RPC và Protocol Buffers}
Giao diện và máy chủ trao đổi dữ liệu qua Connect RPC~\cite{connectrpc} (tương thích gRPC) với hợp đồng được mô tả bằng Protocol Buffers~\cite{protobuf}. Cách tiếp cận này trực tiếp đáp ứng yêu cầu dễ bảo trì ở Chương~\ref{chapter:Related_works}.

\textbf{Hợp đồng tập trung, sinh mã hai phía} -- giao diện trao đổi dữ liệu được mô tả một lần bằng Protocol Buffers rồi sinh mã tự động cho cả Go phía máy chủ và TypeScript phía giao diện, nhờ đó hai phía luôn đồng bộ và giảm sai sót (Hình~\ref{fig:protobuf}).

\begin{figure}[H]
\centering
\includegraphics[width=0.6\textwidth]{Hinhve/chuong_3/protobuf.pdf}
\caption{Sinh mã tự động cho máy chủ và giao diện từ một hợp đồng Protocol Buffers}
\label{fig:protobuf}
\end{figure}

\textbf{An toàn kiểu} -- kiểu dữ liệu được kiểm tra chặt chẽ ngay khi biên dịch, giảm lỗi khi hai phía giao tiếp.

Các lựa chọn thay thế gồm REST kèm JSON (phổ biến nhưng không bảo đảm an toàn kiểu khi biên dịch), gRPC thuần (cần thêm thành phần trung gian để trình duyệt sử dụng) và GraphQL (linh hoạt nhưng làm máy chủ phức tạp hơn). Hệ thống chọn Connect RPC và Protocol Buffers vì vừa an toàn kiểu vừa sinh mã cho cả hai phía từ một định nghĩa chung.

\subsection{Lưu trữ dữ liệu}
Hệ thống lưu trữ dữ liệu trên ba kho khác nhau, mỗi kho phù hợp với một loại dữ liệu, đáp ứng yêu cầu lưu trữ bền vững và an toàn ở Chương~\ref{chapter:Related_works}.

\textbf{PostgreSQL} -- cơ sở dữ liệu quan hệ~\cite{postgresql} lưu các dữ liệu nghiệp vụ chính như người dùng, tổ chức, khóa học và bài giảng.

\textbf{FoundationDB} -- kho khóa--giá trị giao dịch~\cite{foundationdb}, phù hợp lưu các khối dữ liệu lớn như bản chép lời và tiến độ xem của người học.

\textbf{SeaweedFS} -- kho lưu trữ đối tượng tương thích S3\footnote{\url{https://github.com/seaweedfs/seaweedfs}}, lưu video tải lên và âm thanh sinh ra.

Các lựa chọn thay thế gồm MySQL hay MongoDB cho cơ sở dữ liệu; hệ thống chọn cách kết hợp trên để mỗi loại dữ liệu được lưu ở nơi phù hợp nhất với đặc tính của nó.

\subsection{Hàng đợi tác vụ bền vững}
Hệ thống xử lý các tác vụ phân tích video chạy dài một cách bất đồng bộ qua một hàng đợi tác vụ bền vững xây dựng trên PostgreSQL, đáp ứng yêu cầu xử lý ở chế độ nền và phục hồi tác vụ ở Chương~\ref{chapter:Related_works}. Đây cũng là nền tảng cho kiến trúc quản lý các tác vụ phân tích đã giới thiệu ở Chương~\ref{chapter:Introduction}. Hình~\ref{fig:hangdoi} minh họa luồng một tác vụ đi qua hàng đợi.

\begin{figure}[H]
\centering
\includegraphics[width=0.9\textwidth]{Hinhve/chuong_3/hangdoi.pdf}
\caption{Luồng xử lý tác vụ qua hàng đợi bền vững}
\label{fig:hangdoi}
\end{figure}

\textbf{Bền vững} -- mọi tác vụ được lưu trong cơ sở dữ liệu nên không thất lạc khi máy chủ gặp sự cố.

\textbf{Có khả năng phục hồi} -- một tác vụ bị gián đoạn được tự động đưa trở lại hàng đợi để xử lý lại thay vì bị bỏ dở.

Các lựa chọn thay thế gồm hàng đợi trên Redis hoặc hệ môi giới thông điệp như RabbitMQ, Kafka. Bảng~\ref{table:queue} so sánh các lựa chọn theo hạ tầng bổ sung, độ bền và độ phức tạp vận hành.

\begin{table}[ht]
\centering
\small
\caption{So sánh các cơ chế hàng đợi tác vụ}
\label{table:queue}
\begin{tabular}{|>{\raggedright\arraybackslash}p{3.3cm}|>{\raggedright\arraybackslash}p{3cm}|>{\raggedright\arraybackslash}p{2.8cm}|>{\raggedright\arraybackslash}p{3cm}|}
\hline
\textbf{Tiêu chí} & \textbf{PostgreSQL} & \textbf{Redis} & \textbf{RabbitMQ, Kafka} \\ \hline
Hạ tầng bổ sung & Không & Có & Có \\ \hline
Đảm bảo không mất tác vụ & Cao & Trung bình & Cao \\ \hline
Độ bền dữ liệu & Cao & Trung bình & Cao \\ \hline
Thông lượng tối đa & Trung bình & Cao & Rất cao \\ \hline
Độ phức tạp vận hành & Thấp & Trung bình & Cao \\ \hline
\end{tabular}
\end{table}

Như Bảng~\ref{table:queue} cho thấy, hàng đợi trên nền PostgreSQL không đòi hỏi thêm hạ tầng, đổi lại thông lượng tối đa thấp hơn các hệ môi giới chuyên dụng; với quy mô của hệ thống, đây là đánh đổi hợp lý. Cơ chế chi tiết của kiến trúc quản lý tác vụ này được trình bày ở Chương~\ref{chapter:SolutionAndContribution}.

\subsection{Xác thực và bảo mật}
Hệ thống đáp ứng yêu cầu bảo mật ở Chương~\ref{chapter:Related_works} bằng ba cơ chế chính.

\textbf{Xác thực phi trạng thái} -- người dùng được xác thực bằng JSON Web Token~\cite{rfc7519}, giúp máy chủ không cần lưu trạng thái phiên. Hình~\ref{fig:jwt} minh họa quy trình đăng nhập và gọi dịch vụ kèm token. Ở phía giao diện, token được lưu trong cookie HTTP-only và tự động làm mới, nên người dùng không phải đăng nhập lại thường xuyên.

\begin{figure}[H]
\centering
\includegraphics[width=0.85\textwidth]{Hinhve/chuong_3/jwt.pdf}
\caption{Quy trình xác thực bằng JSON Web Token}
\label{fig:jwt}
\end{figure}

\textbf{Bảo vệ mật khẩu} -- mật khẩu người dùng được băm bằng thuật toán bcrypt~\cite{provos1999bcrypt}, vốn có thể điều chỉnh độ phức tạp để thích ứng với sự tiến bộ của phần cứng.

\textbf{Phân quyền theo vai trò} -- quyền truy cập được kiểm soát theo vai trò trên ba phạm vi gồm hệ thống, tổ chức và khóa học, bảo đảm mỗi người dùng chỉ truy cập phần dữ liệu thuộc quyền của mình.

Các lựa chọn thay thế gồm lưu phiên phía máy chủ, dùng giao thức OAuth, hoặc băm mật khẩu bằng Argon2 hay scrypt; hệ thống chọn các phương án trên vì phù hợp với kiến trúc nhiều dịch vụ và không phụ thuộc nhà cung cấp bên ngoài.

\section{Công nghệ phía giao diện}
\label{section:3.4}
Giao diện web là nơi giảng viên quản lý bài giảng và người học tham gia các hoạt động tương tác. Phần này trình bày các công nghệ chính được sử dụng ở phía giao diện.

\subsection{Next.js và React}
Giao diện web được xây dựng bằng thư viện React~\cite{react} trên nền tảng Next.js~\cite{nextjs}, đáp ứng yêu cầu về tính dễ dùng cùng các yêu cầu phi chức năng về hiệu năng và bảo mật ở Chương~\ref{chapter:Related_works}.

\textbf{Kết xuất phía máy chủ} -- Next.js kết xuất trang ở phía máy chủ giúp trang tải nhanh.

\textbf{Xác thực tích hợp} -- Next.js hỗ trợ kiểm soát xác thực ngay ở phía máy chủ trước khi người dùng truy cập các trang cần đăng nhập; Hình~\ref{fig:nextjs} minh họa luồng xác thực và kết xuất trang.

\textbf{Thành phần tái sử dụng} -- React cho phép dựng giao diện từ các thành phần tái sử dụng. Toàn bộ giao diện được trình bày bằng tiếng Việt theo đúng nhóm người dùng mục tiêu.

\begin{figure}[H]
\centering
\includegraphics[width=0.95\textwidth]{Hinhve/chuong_3/nextjs.pdf}
\caption{Luồng xác thực tại tầng proxy và kết xuất trang phía máy chủ của Next.js}
\label{fig:nextjs}
\end{figure}

\textbf{Trình phát video tương tác} -- giao diện học tập sử dụng một trình phát video tùy biến cho phép chèn hoạt động vào đúng mốc thời gian, tạm dừng để người học trả lời, hạn chế tua nhanh qua phần chưa xem và tự động lưu lại tiến độ học; mỗi loại hoạt động được ánh xạ tới một thành phần hiển thị riêng để dễ bổ sung loại mới.

Các lựa chọn thay thế gồm Angular, Vue cùng Nuxt, hay ứng dụng một trang thuần; hệ thống chọn Next.js và React vì hệ sinh thái lớn và hỗ trợ tốt việc kết xuất phía máy chủ.

\subsection{Thư viện giao diện}
Phần trình bày được xây dựng bằng TailwindCSS~\cite{tailwindcss} kết hợp các thành phần giao diện từ Radix và Shadcn\footnote{\url{https://ui.shadcn.com}}, đáp ứng yêu cầu về tính dễ dùng ở Chương~\ref{chapter:Related_works}.

\textbf{Dựng giao diện nhanh và nhất quán} -- TailwindCSS cung cấp bộ tiện ích định kiểu giúp xây dựng giao diện nhanh và đồng nhất về phong cách.

\textbf{Chú trọng khả năng truy cập} -- Radix và Shadcn cung cấp các thành phần chú trọng khả năng truy cập và hiển thị tốt trên nhiều kích thước màn hình.

Các lựa chọn thay thế gồm các bộ thư viện như MUI hay Bootstrap; TailwindCSS được ưu tiên nhờ tính linh hoạt và khả năng tùy biến cao theo thiết kế riêng.

Như vậy, chương này đã trình bày các nền tảng lý thuyết về học tập làm cơ sở cho thiết kế tương tác, cùng các công nghệ cốt lõi ở khối xử lý trí tuệ nhân tạo, phía máy chủ và phía giao diện, kèm lý do lựa chọn so với các phương án thay thế. Việc vận dụng các công nghệ này để phân tích thiết kế, triển khai và đánh giá hệ thống được trình bày trong Chương~\ref{chapter:Experiment}.

\end{document}